\begin{table}[tb]
\centering\footnotesize\setlength{\tabcolsep}{3pt}
\caption{固定した全12条件．$+?$は既知小計に使用量不明の生成が残ることを示す．}
\begin{tabular}{lllrrr}
課題 & 条件 & 成果 & 論理token & 生成 & API \\ \hline
写真 & N & 成功 & 135,628 & 16 & 99 \\
写真 & G & 成功 & 135,628 & 16 & 99 \\
写真 & F & 成功 & 135,628 & 16 & 99 \\
精算 & G & 成功 & 211,110 & 24 & 25 \\
精算 & F & 成功 & 211,110 & 24 & 25 \\
精算 & N & 成功 & 211,110 & 24 & 25 \\
返金 & F & 成功 & 102,087 & 13 & 13 \\
返金 & N & 成功 & 102,087 & 13 & 13 \\
返金 & G & 成功 & 102,087 & 13 & 13 \\
書出し & N & 未達 & 320,456$+?$ & 20 & 42 \\
書出し & G & 未達 & 348,819$+?$ & 21 & 79 \\
書出し & F & 未達 & 423,816 & 23 & 79 \\
\hline\end{tabular}
\end{table}
各条件の成功／評価はNが3/4，Gが3/4，Fが3/4であった．
CSVの最終結果はNが未達，Gが未達，Fが未達であった．
\begin{table}[tb]
\centering\footnotesize\setlength{\tabcolsep}{3pt}
\caption{全４課題を含む３対照．欠測時は既知分だけで全費用差を確定しない．}
\begin{tabular}{llrl}
比較 & 課題別成功 & 全論理token差 & 固定基準 \\ \hline
G−N & 喪失なし & 未確定 & 未達 \\
F−N & 喪失なし & 未確定 & 未達 \\
F−G & 喪失なし & 未確定 & 未達 \\
\hline\end{tabular}
\end{table}
G/Fの通知前の入力と予算が一致し，両側で送信を試行した課題数は1であった．事実通知の便益候補の固定基準は未達だった．これは新手法の優位性や統計的な非劣性の判定ではない．
\begin{table}[tb]
\centering\footnotesize\setlength{\tabcolsep}{3pt}
\caption{失敗を含む費用寄与．全12条件の使用量が確定したときだけ全体差を分解する．}
\begin{tabular}{llrr}
比較 & 成否の組合せ & 課題数 & 論理token差 \\ \hline
G−N & 両成功 & 3 & 未確定 \\
G−N & 両未達 & 1 & 未確定 \\
G−N & 成否差 & 0 & 未確定 \\
G−N & 未評価あり & 0 & 未確定 \\
F−N & 両成功 & 3 & 未確定 \\
F−N & 両未達 & 1 & 未確定 \\
F−N & 成否差 & 0 & 未確定 \\
F−N & 未評価あり & 0 & 未確定 \\
F−G & 両成功 & 3 & 未確定 \\
F−G & 両未達 & 1 & 未確定 \\
F−G & 成否差 & 0 & 未確定 \\
F−G & 未評価あり & 0 & 未確定 \\
\hline\end{tabular}
\end{table}
未達を安く終えた差を，完了した仕事の効率改善に読み替えない．未実行・未評価を未達で補完せず，欠測は残した．
